\chapter{Ventaja y disrupción}\label{ch:advantage-and-disruption}

% Alcance: visión basada en recursos, VRIN, Océano Azul, el dilema del innovador.
% Enlace cruzado: la ventaja durable es lo que sostiene LTV > CAC a lo largo del tiempo (\cref{ch:customer-economics}).

El análisis estructural (\cref{ch:competitive-analysis}) explica los retornos a nivel de industria;
no puede explicar por qué algunas empresas obtienen retornos por encima del promedio \emph{dentro}
de industrias poco atractivas durante décadas. La visión basada en recursos, la Estrategia de
Océano Azul y el dilema del innovador ofrecen cada una una respuesta diferente a esa pregunta
--- y cada una implica una respuesta estratégica distinta cuando la respuesta cambia.

\section{La visión basada en recursos y VRIN}\label{sec:rbv-vrin}
% complejidad: 3

¿Cuáles de los activos y capacidades de la empresa sostendrán una ventaja competitiva, y cuáles
serán erosionados por la competencia? La auditoría debe preceder a cualquier apuesta de crecimiento.
Las industrias compiten por clientes; las empresas compiten por recursos. Si los recursos que
impulsan el desempeño pueden copiarse o comprarse, los retornos revierten al costo de capital
independientemente de la cuota de mercado. La ventaja durable requiere recursos que, en un
sentido técnico preciso, no sean replicables por los rivales.

\textcite{barney1991resource} define un recurso como fuente de ventaja competitiva sostenida
únicamente cuando satisface los cuatro criterios de manera simultánea:
\begin{description}
  \item[Valioso] El recurso permite a la empresa explotar oportunidades o neutralizar amenazas
        en su entorno --- es decir, está vinculado a la disposición a pagar del cliente o a la
        reducción de costos. Un recurso que no supera ninguna prueba de las Cinco Fuerzas
        (\cref{ch:competitive-analysis}) no es valioso en el sentido de Barney.
  \item[Raro] El recurso no está en posesión de un gran número de competidores actuales o
        potenciales. Los recursos comunes generan paridad competitiva, no ventaja.
  \item[Inimitable] Los rivales no pueden copiar ni adquirir fácilmente el recurso. La
        inimitabilidad surge de la dependencia de la trayectoria (acumulada mediante una historia
        irrepetible), la ambigüedad causal (los competidores no pueden identificar qué recurso
        causa el desempeño) o la complejidad social (por ejemplo, la cultura organizacional,
        las alianzas basadas en la confianza).
  \item[No sustituible] No existe ningún recurso alternativo estratégicamente equivalente.
        Incluso un recurso raro e inimitable pierde valor si un sustituto entrega el mismo
        resultado estratégico.
\end{description}
Estos cuatro criterios forman un árbol de decisión, no una lista de verificación: si se falla
en cualquiera de ellos, el recurso obtiene a lo sumo una ventaja transitoria. Los recursos que
superan los cuatro generan \emph{rentas ricardianas} --- retornos por encima del costo de
oportunidad del recurso que persisten porque la oferta del recurso es inelástica.

VRIN es diagnóstico, no prescriptivo: identifica dónde vive actualmente la ventaja, pero no
ofrece una teoría sobre cómo construir nuevos recursos. El \emph{framework} también subestima
la importancia de las \emph{capacidades dinámicas} --- la habilidad de la empresa para
reconfigurar su base de recursos conforme el entorno cambia. Un recurso que hoy es VRIN puede
quedar obsoleto por una discontinuidad tecnológica (\cref{sec:innovators-dilemma}), momento en
el que la auditoría estática de Barney devuelve la respuesta equivocada.

\begin{example}
\emph{Decisión: ¿en qué debería invertir la empresa SaaS para proteger su margen?}

Auditoría de recursos candidatos frente a VRIN:

\begin{itemize}
  \item \textbf{Datos propietarios de clientes (comportamiento + registros de integración)}: Valioso
        (permite personalización e ingresos por expansión); Raro (el conjunto de datos de cada
        empresa es único); Inimitable (acumulación dependiente de la trayectoria --- no puede
        comprarse al por mayor); No sustituible (los datos de intención de terceros no replican
        las señales longitudinales de uso). \emph{Veredicto: VRIN --- el núcleo defendible.}
  \item \textbf{Ecosistema de integraciones (más de 50 conectores nativos)}: Valioso y Raro en el
        lanzamiento; Inimitable si la profundidad de integración supera un umbral de replicación;
        No sustituible si el costo de cambio (\cref{ch:growth}) es suficientemente alto.
        \emph{Veredicto: parcialmente VRIN --- invertir para profundizar.}
  \item \textbf{Nombre de marca}: Valioso (reduce el CAC mediante búsqueda orgánica); aún no
        Raro en un mercado fragmentado. \emph{Veredicto: no VRIN --- apoya, no protege.}
  \item \textbf{Base de código}: No es raro (los \emph{frameworks} modernos están estandarizados)
        e imitable con suficiente gasto en ingeniería. \emph{Veredicto: paridad competitiva, no
        ventaja.}
\end{itemize}

La inversión en I+D debería profundizar la infraestructura de datos y las integraciones
(recursos que superan el criterio VRIN) en lugar de reconstruir funcionalidades genéricas que
los rivales ya ofrecen. Esto eleva directamente el LTV mediante menor \emph{churn}
(\cref{ch:customer-economics}) y es el mecanismo por el cual \cref{eq:roic} supera al WACC
de manera sostenida.
\end{example}

\begin{admonition}{Advertencia}
Clasificar todo lo que tiene la empresa como «recurso estratégico». Los criterios VRIN son
restrictivos por diseño. La mayoría de los recursos son \emph{necesidades competitivas} ---
requeridas para competir, pero que no confieren ninguna ventaja. Confundir una necesidad (un
CRM) con una ventaja (los datos propietarios que contiene) desvía el capital hacia mantener
la paridad en lugar de construir rentas.
\end{admonition}

\textcite{dranove7theconomics} sintetizan la visión basada en recursos con la organización
industrial: las rentas ricardianas requieren tanto recursos VRIN \emph{como} ciertas barreras
a la competencia en el mercado de recursos. Un recurso no puede ser VRIN en el mercado de
productos si se negocia libremente en el mercado de recursos a su NPV completo --- los
mercados de recursos competitivos disipan las rentas antes de que el mercado de productos las
perciba.

La auditoría VRIN completa el panorama estructural que comenzó con las Cinco Fuerzas:
\cref{ch:competitive-analysis} identifica las amenazas externas; VRIN identifica los recursos
que las resisten. El ROIC sostenido que supera al WACC (\cref{eq:advantage-spread} más
adelante) es la expresión financiera de un recurso que supera esta prueba. \textcite{barney1991resource}
es la fuente primaria; \textcite{dranove7theconomics} aportan la integración con la organización
industrial.

\section{Estrategia de Océano Azul}\label{sec:blue-ocean}
% complejidad: 3

En lugar de competir con más fuerza en un mercado existente en disputa, ¿se puede definir una
nueva curva de valor que haga a la competencia temporalmente irrelevante? La mayoría de la
estrategia trata de ganar una mayor parte de un pastel existente. El Océano Azul pregunta si
el pastel puede reformarse reduciendo costos \emph{y} elevando el valor para el comprador
simultáneamente --- eliminando características en las que la industria compite pero que los
clientes no valoran, y creando características que los clientes desean pero que nadie ofrece.

\textcite{kim2015} proponen la cuadrícula ERRC como herramienta operativa:
\begin{description}
  \item[Eliminar] ¿Cuáles de los factores que la industria da por sentados deben eliminarse por
        completo? Estos añaden costo sin añadir valor percibido; eliminarlos financia la columna
        «Crear».
  \item[Reducir] ¿Cuáles factores deben reducirse muy por debajo del estándar de la industria?
        Los competidores sobre-entregan en estos; el ahorro de costo es redistribuible.
  \item[Elevar] ¿Cuáles factores deben elevarse muy por encima del estándar de la industria?
        Estos abordan necesidades genuinas no satisfechas.
  \item[Crear] ¿Qué factores deben introducirse que la industria nunca ha ofrecido? Este es el
        mecanismo de creación de demanda; apunta al no-consumo en lugar de captar clientes de
        los rivales.
\end{description}
El resultado es una \emph{curva de valor} --- un gráfico del nivel de oferta a lo largo de los
factores competitivos --- que diverge visiblemente del promedio de la industria. La divergencia
es la prueba: una curva que simplemente puntúa más alto en todo no es un movimiento de Océano
Azul, es una estrategia premium.

La cuadrícula ERRC es direccional y cualitativa; no proporciona ningún mecanismo para estimar
si la nueva curva de valor genera suficiente demanda para cubrir la inversión. Una oferta
rediseñada que ningún segmento de clientes valora supera la prueba ERRC pero fracasa
comercialmente. Combinar con la perspectiva de segmentación de \cref{ch:segmentation-targeting-positioning}
para verificar a qué no-consumidores llega realmente el movimiento.

\begin{example}
\emph{Decisión: ¿puede la empresa SaaS crear un Océano Azul frente a Salesforce y CRM
empresariales similares?}

\begin{itemize}
  \item \textbf{Eliminar}: Implementación compleja (despliegues de 6 a 12 meses), roles de
        administrador dedicados, opciones de datos en las instalaciones.
  \item \textbf{Reducir}: Precio en comparación con el CRM empresarial (Salesforce tiene un
        precio de \money{150}--\money{300}/usuario/mes; nuestro \money{50}/puesto es un
        recorte deliberado), conjunto de funcionalidades en relación con el Salesforce Sales
        Cloud completo.
  \item \textbf{Elevar}: Tiempo hasta el primer valor (objetivo: en funcionamiento en un día),
        integraciones nativas con herramientas modernas de productividad, facturación por uso
        transparente.
  \item \textbf{Crear}: \emph{Onboarding} de autoservicio sin necesidad de contacto comercial;
        flujos de trabajo guiados dentro del producto diseñados para operadores no técnicos.
\end{itemize}

La curva apunta a \emph{no-consumidores} del CRM --- pymes que no podían absorber el costo de
implementación empresarial --- en lugar de los clientes existentes de Salesforce. Si tiene
éxito, el CAC de este segmento es menor (no se requiere equipo de ventas para el autoservicio)
y el \emph{churn} es menor (mayor velocidad hacia el primer valor). Ambos mejoran la razón
LTV:CAC en \cref{ch:customer-economics}.
\end{example}

\begin{admonition}{Advertencia}
Aplicar la cuadrícula ERRC de manera \emph{a posteriori} para racionalizar el posicionamiento
existente en lugar de usarla como \emph{herramienta de diseño prospectivo}. Los críticos
(\parencite{dranove7theconomics}) señalan que la mayoría de los casos de Océano Azul se
escriben después del hecho, cuando la curva de valor que resultó exitosa recibe la etiqueta
de un movimiento ERRC deliberado. Usada rigurosamente, la cuadrícula es una hipótesis sobre
la demanda; validarla con experimentos de mercado antes de comprometer capital.
\end{admonition}

\begin{admonition}{Controvertido}
La Estrategia de Océano Azul es el \emph{framework} estratégico más controvertido de este
libro. La objeción central de los críticos: los casos exitosos de Océano Azul sólo son
distinguibles de apuestas de producto afortunadas en retrospectiva. La cuadrícula ERRC no
indica \emph{qué} factores eliminar o crear --- eso requiere una perspectiva del lado de la
demanda que el \emph{framework} no suministra. Los defensores responden que la visualización
de la curva de valor disciplina a un equipo directivo para alejarse del pensamiento
incremental. Ambos lados tienen razón; tratar el Océano Azul como un estímulo a la
creatividad validado por evidencia del cliente, no como una metodología de planificación.
\end{admonition}

La Estrategia de Océano Azul se intersecta con la visión de la demanda desde los
trabajos-por-realizar (\cref{ch:customer-economics}): la columna «Crear» tiene éxito cuando
aborda un trabajo que ningún producto existente realiza bien. El mecanismo está más cerca de
la innovación disruptiva (más adelante) que de la diferenciación tradicional, porque apunta a
criterios de éxito distintos en lugar de superar a los incumbentes en criterios compartidos.

La cuadrícula ERRC responde la pregunta de diseño; la auditoría VRIN (\cref{sec:rbv-vrin})
pregunta luego si la oferta resultante es defendible una vez que los rivales pueden observarla.
Un Océano Azul que supera el criterio VRIN es duradero; uno que no lo hace será imitado y la
ventaja se comprimirá. \textcite{kim2015} es la fuente primaria.

\section{El dilema del innovador}\label{sec:innovators-dilemma}
% complejidad: 5 -> figure disruption-scurve

¿Es un entrante aparentemente inferior una amenaza de disrupción genuina, o simplemente un
producto diferente? La respuesta determina si se debe responder agresivamente, ignorar o
adquirir. Los incumbentes pierden no porque estén mal gestionados sino porque están bien
gestionados: escuchan a sus mejores clientes, invierten donde los márgenes son más altos y
mejoran el desempeño en las dimensiones que sus clientes más valoran. Los disruptores entran
en dimensiones que ningún cliente actual solicita, sirven a no-consumidores o a los
sub-atendidos, y mejoran hasta superar al segmento \emph{mainstream}. El peligro es invisible
hasta que es demasiado tarde.

\textcite{christensen1997} identifica dos tipos de innovación:
\begin{description}
  \item[Innovaciones incrementales] Mejoran el desempeño en las dimensiones que los clientes
        existentes valoran. Los incumbentes casi siempre ganan estos concursos; tienen escala,
        canales establecidos e incentivos correctos.
  \item[Innovaciones disruptivas] Introducen una dimensión de desempeño diferente --- a menudo
        inicialmente inferior en la métrica establecida (desempeño, funcionalidades,
        confiabilidad) pero superior en accesibilidad, precio o simplicidad. Entran en segmentos
        de no-consumo o del extremo inferior que el incumbente no defiende.
\end{description}
El mecanismo de la disrupción es un desfase de curvas en S: la trayectoria incremental del
incumbente supera el techo de demanda del mercado \emph{mainstream}. La trayectoria de mejora
del disruptor comienza por debajo del umbral pero, a medida que asciende en la curva en S de
su propia tecnología, eventualmente cruza el umbral de desempeño \emph{mainstream} --- punto
en el que la base de clientes del incumbente migra en masa porque el disruptor es ahora
«suficientemente bueno» y más barato.

\begin{figure}[htbp]
  \centering
  \includegraphics[width=0.86\linewidth]{disruption-scurve}
  \caption{El dilema del innovador como curvas en S duales. La trayectoria del incumbente supera
    el umbral de demanda \emph{mainstream}; el disruptor entra por debajo de él en un segmento
    sub-atendido y cruza el umbral a medida que su curva en S madura. El punto de cruce es
    cuando ocurre la disrupción.}
  \label{fig:disruption-scurve}
\end{figure}

El criterio de Christensen para la disrupción es la entrada en segmentos \emph{sub-atendidos
o de no-consumo} --- no simplemente «entrada desde abajo» en sentido de precio. Todo incumbente
llama «disruptivo» a todo entrante de bajo precio; el poder predictivo del \emph{framework}
depende de identificar correctamente si la trayectoria de desempeño del entrante eventualmente
satisfará la demanda \emph{mainstream}. Muchos entrantes del extremo inferior permanecen
confinados en segmentos del extremo inferior de forma permanente (invasión del extremo inferior,
no disrupción). La prueba es la trayectoria, no el precio actual.

\begin{example}
\emph{Decisión: ¿debería la empresa SaaS tratar a las herramientas de automatización sin
código (Zapier, Make, Airtable) como una amenaza disruptiva?}

Las herramientas sin código actualmente puntúan bajo en las dimensiones que los clientes
existentes de la SaaS valoran (profundidad de datos, reportes, cumplimiento normativo). No
son competitivas \emph{hoy} para el comprador del mercado medio.

La pregunta disruptiva es: ¿está la curva en S del sin-código mejorando con suficiente rapidez
como para alcanzar los requisitos de desempeño del mercado medio dentro del horizonte de
planificación? La evidencia sugiere que sí: el financiamiento de Airtable en 2022 apuntó
explícitamente a los flujos de trabajo del mercado medio. La prueba de Christensen es si las
herramientas sin código están mejorando en la \emph{trayectoria correcta}, no si son
competitivas hoy.

Respuesta: invertir en los recursos VRIN (profundidad de datos, integraciones) que las
herramientas sin código no pueden replicar rápidamente, mientras se monitorea la tasa de
mejora del sin-código como indicador adelantado. Un rezago de dos años entre «entra al
no-consumo» y «cruza el umbral \emph{mainstream}» es históricamente común --- esa es la
ventana de respuesta.
\end{example}

\begin{admonition}{Advertencia}
Todo incumbente etiqueta a todo entrante como «disruptivo» porque la palabra se ha convertido
en un cumplido. \emph{Disruptivo} en el sentido de Christensen es un término técnico preciso:
significa que el entrante compite inicialmente en criterios de desempeño diferentes en un
segmento de mercado diferente. Un entrante premium que apunta a los mejores clientes del
incumbente es una amenaza \emph{incremental} (más difícil de superar en funcionalidades, no
disrupción). Clasificar erróneamente el tipo de amenaza produce la respuesta incorrecta: la
disrupción requiere una unidad de negocio de extremo inferior separada; las amenazas
incrementales requieren inversión directa en el producto.
\end{admonition}

El dilema es tanto organizacional como tecnológico: un incumbente racional asigna recursos
hacia sus clientes de mayor margen, precisamente los que el disruptor no apunta inicialmente.
La solución de Christensen --- una unidad estructuralmente separada con estructura de costos y
métricas diferentes --- es un problema de gobierno corporativo (\cref{ch:organizational-behavior}).
La expresión financiera del dilema es un ROIC que luce excelente justo hasta el cruce de la
disrupción, y luego colapsa (\cref{eq:advantage-spread}).

La disrupción reinicia la auditoría VRIN: un recurso previamente inimitable (la distribución
física, en el caso de la música) puede ser sustituido por una nueva infraestructura tecnológica.
La sección final del capítulo vincula la dinámica de la disrupción con la métrica financiera
que la hace legible. \textcite{christensen1997} es la fuente primaria.

\section{La ventaja como diferencial de valor}\label{sec:advantage-spread}
% complejidad: 3

¿Cómo sabemos si la empresa posee actualmente una ventaja competitiva, y si esta se está
ampliando o estrechando con el tiempo? Los \emph{frameworks} cualitativos diagnostican la
\emph{fuente} de la ventaja; se necesita una métrica financiera para confirmar su
\emph{magnitud} y rastrear si la estrategia está funcionando. La medida natural es si los
retornos superan el costo de capital por un margen suficiente para justificar el riesgo ---
la misma prueba que aplicaría un accionista.

La ventaja competitiva se define operacionalmente como un diferencial de valor positivo
sostenido:
\begin{equation}\label{eq:advantage-spread}
  \text{Diferencial de ventaja} \;=\; \mathrm{ROIC} - \mathrm{WACC},
\end{equation}
donde ROIC se define en \cref{eq:roic} y WACC en \cref{eq:wacc}. Cuando el diferencial es
positivo, la empresa obtiene rentas ricardianas y crea beneficio económico
(\cref{eq:economic-profit}). Cuando es cero o negativo, la empresa obtiene paridad
competitiva o destruye valor independientemente del beneficio contable.

El diferencial se conecta directamente con la economía del cliente: un ROIC sostenido
$>$ WACC requiere un LTV sostenido $>$ CAC (\cref{eq:ltv-cac} en \cref{ch:customer-economics}),
porque el valor presente neto de cada cohorte debe ser positivo y creciente. Un diferencial de
ventaja en expansión es la expresión agregada de un \emph{moat} en profundización.

El ROIC puede manipularse mediante elecciones contables (gastos vs. capitalización de I+D,
pasivos fuera del balance). Una empresa que capitaliza el desarrollo de software infla la base
de capital invertido y suprime el ROIC; una que lo gasta deflacta la base. Normalizar antes de
comparar entre empresas o a lo largo del tiempo. Del mismo modo, un diferencial positivo en un
solo año puede reflejar un pico cíclico en lugar de una ventaja estructural --- el diferencial
debe ser persistente para ser evidencia de un \emph{moat}.

\begin{example}
\emph{Decisión: ¿tiene actualmente la empresa SaaS una ventaja medible?}

Según la hoja de datos: NOPAT $\approx$\money{600000}, capital invertido $\approx$\money{3000000},
lo que da ROIC $= \money{600000}/\money{3000000} = \qty{20}{\percent}$. WACC $= \qty{11}{\percent}$
(\cref{eq:wacc}).
\[
  \text{Diferencial de ventaja} = \qty{20}{\percent} - \qty{11}{\percent} = \qty{9}{\percent}.
\]
El diferencial es positivo y material. Pero debe evaluarse frente al análisis del \emph{moat}:
si la auditoría VRIN muestra sólo ventajas transitorias, el diferencial se comprimirá a medida
que los rivales repliquen. Un diferencial del \qty{9}{\percent} sostenido durante cinco años
con \money{3000000} de capital invertido crea $5 \times 0{,}09 \times \money{3000000} = \money{1350000}$
de beneficio económico acumulado --- la recompensa por recursos VRIN genuinos.

Si la inversión en el \emph{moat} (\cref{ch:growth}) eleva el LTV en \money{500} (por ejemplo,
menor \emph{churn} gracias a integraciones más profundas), la empresa puede gastar \money{500}
más por cliente adquirido sin violar el piso de LTV:CAC (\cref{eq:ltv-cac}).
\end{example}

\begin{admonition}{Advertencia}
Comparar el ROIC con el costo contable de la deuda en lugar de con el WACC. La deuda es barata
en términos nominales; el WACC incorpora el retorno requerido sobre el \emph{equity}, que es
considerablemente mayor. Una empresa que obtiene \qty{8}{\percent} de ROIC y se financia a
\qty{5}{\percent} de deuda nominal \emph{parece} rentable pero destruye valor de \emph{equity}
si el WACC $=\qty{11}{\percent}$.
\end{admonition}

El \emph{framework} de valoración de McKinsey (\parencite{dranove7theconomics}) muestra que el
valor de mercado es el valor presente de los futuros beneficios económicos descontados, no de
los beneficios contables. Una empresa que cotiza con prima respecto al valor en libros tiene
un NPV positivo del diferencial de ventaja futuro; una empresa que cotiza con descuento indica
que el mercado espera que el ROIC caiga por debajo del WACC. Los múltiplos de EV en
\cref{ch:valuation} son el pronóstico del mercado sobre el diferencial futuro, no una
instantánea del diferencial actual.

El diferencial de ventaja (\cref{eq:advantage-spread}) vincula los \emph{frameworks} cualitativos
de este capítulo --- VRIN, Océano Azul, disrupción --- con el lenguaje financiero de
\cref{ch:corporate-finance-core} y \cref{ch:valuation}. El diferencial sostenido es la prueba
empírica que distingue la ventaja durable de la buena fortuna temporal. \cref{ch:growth}
pregunta luego cómo ampliarlo mediante la construcción del \emph{moat}.
